% SEGMENT OLP-0386-S01
% Part: many-valued-logic
% Chapter: syntax-and-semantics
% Section: connectives

\documentclass[../../../include/open-logic-section]{subfiles}

\begin{document}

\olfileid{mvl}{syn}{con}

\olsection{மொழிகளும் இணைப்பிகளும்}

செவ்வியல் கூற்றுத் தருக்கமும் வேறு பல தருக்கங்களும் \emph{கூற்று
மாறிலிகள்}, \emph{இணைப்பிகள்} ஆகியவற்றின் ஒரு கணத்தைப் பயன்படுத்துகின்றன.
எடுத்துக்காட்டாக,
பின்வருவனவற்றை மூலக் குறிகளாகப் பயன்படுத்துகிறோம்:
\begin{enumerate}
\tagitem{prvFalse}{!!{falsity}க்கான கூற்று மாறிலி~$\lfalse$.}{}
\tagitem{prvTrue}{!!{truth}க்கான கூற்று மாறிலி~$\ltrue$.}{}
\item தருக்க இணைப்பிகள்:
  \startycommalist
  \iftag{prvNot}{\ycomma $\lnot$ (மறுப்பு)}{}%
  \iftag{prvAnd}{\ycomma $\land$ (இணையல்)}{}%
  \iftag{prvOr}{\ycomma $\lor$ (பிரிப்பிணைவு)}{}%
  \iftag{prvIf}{\ycomma $\lif$ (!!{conditional})}{}%
  \iftag{prvIff}{\ycomma $\liff$ (!!{biconditional})}{}%
\end{enumerate}
\iftag{defNot,defOr,defAnd,defIf,defIff,defTrue,defFalse,defEx,defAll}{%
மேலுள்ள மூல இணைப்பிகளுடன், பின்வருபவை போன்ற சுருக்கக்குறிகளாக
வரையறுக்கப்பட்ட குறிகளையும் பயன்படுத்துகிறோம்:
\startycommalist
  \iftag{defNot}{\ycomma $\lnot$ (மறுப்பு)}{}%
  \iftag{defAnd}{\ycomma $\land$ (இணையல்)}{}%
  \iftag{defOr}{\ycomma $\lor$ (பிரிப்பிணைவு)}{}%
  \iftag{defIf}{\ycomma $\lif$ (!!{conditional})}{}%
  \iftag{defIff}{\ycomma $\liff$ (!!{biconditional})}{}%
  \iftag{defFalse}{\ycomma $\lfalse$ (!!{falsity})}{}%
  \iftag{defTrue}{\ycomma $\ltrue$ (!!{truth})}.}{}

இதே இணைப்பிகள் பல்மதிப்புத் தருக்கங்களிலும் பயன்படுத்தப்படுகின்றன. ஆனால்,
எடுத்துக்காட்டாக இணையலின் வெவ்வேறு வடிவங்களை ஒரே தருக்கத்தில் சேர்ப்பது
பல நேரங்களில் பயனுள்ளது; அந்த வடிவங்களை வேறுபடுத்த வெவ்வேறு குறிகள் தேவை.
சில பல்மதிப்புத் தருக்கங்களில் செவ்வியல் தருக்கத்தில் சமானம் இல்லாத
இணைப்பிகளும் உள்ளன. ஆகவே வழக்கத்தைவிடச் சிறிது பொதுவாகக் கூறுவோம்.

\begin{defn}
  ஒரு \emph{கூற்று மொழி} என்பது \emph{இணைப்பிகளின்} கணம்~$\Lang L$ ஐக்
  கொண்டது. ஒவ்வோர் இணைப்பி~$\star$ க்கும் ஓர் \emph{இடவெண்} உண்டு;
  இடவெண்~$n$ கொண்ட இணைப்பியை \emph{$n$-இட இணைப்பி} என்போம்.
  இடவெண்~$0$ கொண்ட இணைப்பிகளை \emph{மாறிலிகள்} என்றும், இடவெண்~$1$
  கொண்ட இணைப்பிகளை \emph{ஓரிட இணைப்பிகள்} என்றும், இடவெண்~$2$ கொண்ட
  இணைப்பிகளை \emph{ஈரிட இணைப்பிகள்} என்றும் கூறுகிறோம்.
\end{defn}

\begin{ex}
  கூற்றுத் தருக்கத்தின் செந்தர மொழி~$\Lang L_0$ பின்வரும் இணைப்பிகளைக்
  (அவற்றுக்குரிய இடவெண்களுடன்) கொண்டது:
  $\lfalse$~($0$)
  $\lnot$~($1$),
  $\land$~($2$),
  $\lor$~($2$),
  $\lif$~($2$). நாம் கருதும் பெரும்பாலான தருக்கங்கள் இம்மொழியையே
  பயன்படுத்தும்.
  % SOURCE CORRECTION TA-MVL-002: repair “by tradition an convention” while retaining the source's intended symbol contrast.
  மரபாலும் நடைமுறையாலும் சில தருக்கங்கள் சில இணைப்பிகளுக்கு வெவ்வேறு
  குறிகளைப் பயன்படுத்துகின்றன. எடுத்துக்காட்டாக, பெருக்கல் தருக்கத்தில்
  இணையலுக்கான குறி~$\land$ என்பதற்குப் பதிலாகப் பெரும்பாலும்~$\odot$ ஆகும்.
  சில நேரங்களில் தீர்மானத்தன்மைச் செயலி~$\triangle$ ($1$-இடம்) போன்ற
  புதிய செயலியைச் சேர்ப்பது வசதியாக இருக்கும்.
\end{ex}

\end{document}
